\section{What did not work}\label{sec:rejected}
Recorded so that the same ground is not covered twice on this class of machine. The numbers in this list are
measurements from the project log (not in \code{results/}) unless a result file is named. The SMA launch is on its
1.071\,M-cell mesh in the project-log numbers and on its 0.871\,M-cell mesh in the result files of 2026-11-04 and
2026-11-05 (Table~\ref{tab:models}).
\begin{itemize}
\item \textbf{Two simulations side by side} for parameter studies. GCPW alone at 4 threads 1\,785\,\MC; two concurrent
  4-thread runs 967 + 957; two 5-thread runs 864 + 857; SMA one 8-thread run 1\,397 versus two 4-thread runs 691 + 692.
  The bus is shared; pairs give $-4$ to +8\,\%, inside the noise. Studies stay sequential.
\item \textbf{Coarser fused blocks} (\code{OPENEMS\_FUSED\_BLOCK=k} x-lines per block): GCPW +4\,\% (noise), notch
  tutorial 0 at 4 threads and $-7$\,\% at 8, SMA $-14$\,\% (block 4) and $-26$\,\% (block 16) at 4 threads. The L2
  reuse of the line behind is lost with wider blocks. The knob stays at 1.
\item \textbf{CPU field snapshots} so dumps could be evaluated in the background: $-3$\,\%; the 12 small nf2ff dumps
  evaluated concurrently: 0\,\%.
\item \textbf{Two or more time steps per sweep} ($k$-step temporal blocking) on the CPU: +10--14\,\% for $k=2$ in the
  standalone kernel of 2026-10-28 (\code{scripts/kernel/kb\_tb.cpp}: the stock arithmetic with its invariants hoisted
  and the 12 coefficient tables, a fused sweep, $k=1$, against $k=2$, on one independent sub-grid per thread), while
  every extension hook would need rewriting; not pursued then.
  Measured again for cg29 as an upper bound (standalone kernel \code{kb\_tbk}, not in \code{scripts/}: the fork's
  fused IJKN loops, one independent sub-grid per thread, so no halo exchange and no extension hooks, bit-identical for
  every $k$; 6 rounds, ratio to $k=1$): on a 7.9\,M-cell grid 1.038, 1.032 and 0.958$\times$ for $k=2$, 3 and 4 at 4
  threads and 0.990--1.020$\times$ at 8 (0.5 cores of other load); on a 5.25\,M-cell grid 0.971--0.976$\times$ at 4
  threads and 1.055--1.059$\times$ at 8 (1.2 cores). At most +6\,\%, against the +10--14\,\% above; why the fork's
  loops gain less was not examined. A real version (halo lines per thread, every extension's line hooks per sub-step)
  would gain less. Not built (\code{results/cpu\_items\_cg29\_2026-11-05.txt}, section 2). Deeper wavefront or diamond blocking on the
  CPU~\cite{malas2015,fukaya2018} was not tried; on the GPU, diamond blocking is the fork's Metal engine
  (Section~\ref{sec:gpu}, with the tile schedule of luikore's engine, PR \#228~\cite{luikore228}). Both tests gave
  every thread a sub-grid of its own. Upstream master's \code{cpu} engine blocks with all threads in one cache-sized
  tile, which is worth 1.15--1.49$\times$ to it on this machine (Section~\ref{sec:threebuilds}); that form was not
  tried and is open (Section~\ref{sec:future}).
\item \textbf{Compiler flags}: \code{-mcpu=apple-m4} gives an identical instruction stream; link-time optimisation (LTO)
  was not tried; \code{-ffast-math} is ruled out (not bit-identical); GCC~16, measured at one thread in the standalone
  kernel only, runs the stock inner loop 1.16--1.39$\times$ faster than Apple clang but the rewritten one at about the
  same speed (and not bit-identical; \code{results/kernel\_gcc\_vs\_clang\_2026-11-03.txt}, the inner-loop item of
  Section~\ref{sec:engine}). Flush-to-zero: no speed effect on Apple cores (Section~\ref{sec:ceiling}).
\item \textbf{fp16 field storage}: needs field scaling (late-run currents are $10^{-6}$\,A against a half-precision
  normal range ending at $6\times10^{-5}$), touches the 33 Metal kernels and the host mirror of the base's GPU engine,
  or, in the Metal diamond engine, its one kernel and the CPU code that reads the same field arrays (probes, dumps,
  the CPU fallback), and cannot be bit-identical. Not started.
\item \textbf{A Swift rewrite of the engine}: the prototype runs at the cell rate of upstream's Metal engine and of the
  CPU engine of the time (1\,670 versus 1\,638 versus 1\,562\,\MC\ on $201^3$, the CPU on snapshot cg9). Its diamond
  mode (\code{--diamond K}, 2026-11-04) reached 2\,445.7\,\MC\ at $K=4$ on $201^3$ against 1\,540.0 unblocked (one run
  each on a loaded machine, \code{results/diamond\_swift\_2026-11-04.txt}) and became the prototype of the Metal
  diamond engine (Section~\ref{sec:gpu}). The rewrite itself stays a side project.
\item \textbf{The default thread count by wall-clock ramp}: replaced, not tuned.
\item \textbf{The second review's CPU items that gained nothing} (measured for cg29; paired A/B in one build, median
  ratio; \code{results/cpu\_items\_cg29\_2026-11-05.txt}, section 2; item numbers of Table~\ref{tab:review2}): a cache
  of the automatic thread count per model (item 3; at real run lengths GCPW 1.021$\times$ [0.893--1.044], SMA
  1.002$\times$, 5 rounds, 0.5--0.8 cores: the search costs about 1.5\,s of a 10--80\,s run and finds the same count);
  a per-thread QoS split between P and E cores, a 200\,$\mu$s spin and 5\,ms slab balancing (item 4; at 8 threads on
  five models 0.949--1.357$\times$, 0.907--1.122$\times$ and 0.991--1.008$\times$, 5 rounds, 0.3--0.5 cores; with the
  flags the QoS split lost on three of six models at the automatic count, 0.886--0.948$\times$: gains on two or three
  models, losses of 3--11\,\% on the others); Mur z-face vectors (item 6; 0.982--1.012$\times$, 6 rounds, 0.4--2
  cores); per-line hook masks (item 7; 0.960--1.020$\times$, GCPW 1.005$\times$ and 1.009$\times$ at 4 and 8 threads,
  5--8 rounds); the split form of item 9, the distinct vectors in one table and four 16-bit ids per operator (SMA
  0.952--0.999$\times$, GCPW 0.933--0.954$\times$, patch x5 and notch x3 0.967--1.023$\times$, 5 rounds: the extra
  loads cost more than the smaller tables save; the geometry tables of Section~\ref{sec:cg29} replace it); and a
  50\,ms tuning window (short runs GCPW 1.053$\times$, notch tutorial 1.100$\times$, but at real run lengths GCPW
  1.017$\times$ [0.861--1.120] and SMA 0.993$\times$; the default stays 150\,ms). Not built: item 10, spin barriers
  for the start and stop of each processing interval (the boards' calls last 50--300\,ms, so the two wake-ups of
  20--40\,$\mu$s per call are under 0.1\,\%; the patch tutorial could gain 1--2\,\% at most, an estimate). Skipped:
  item 8, since the Metal diamond engine replaces the base's GPU engine. Of other contributors' ideas, two were already
  in the fork, MrHighVoltage's thread tuner that scores real time steps (since cg8) and PR \#231~\cite{akdag231} (the
  hook plan and the UPML rows along z), and two parts were not needed: the Debye half of PR \#244 (the Debye slab's
  setup takes 0.02\,s) and the threaded unfused dispersive hooks of PR \#245.
\item \textbf{On the GPU} (the Metal diamond engine, Section~\ref{sec:gpu}): a fused but unblocked sweep ($K=1$) runs
  the two large models at 0.57--0.65$\times$ of the default $K=4$ (0.54$\times$ on patch x5 at the widest tile; paired,
  4 rounds, 2.5--4 cores of other load; \code{results/fused\_sweep\_k1\_2026-11-04.txt}); the UPML folded into the cell
  update gained nothing (paired against cg24, 4 rounds: SMA 1.023$\times$ [0.934--1.106], notch tutorial 0.948$\times$,
  notch x3 0.988$\times$; \code{results/diamond\_engine\_cg24\_cg27\_2026-11-04.txt}); the I-J-K-N field order of the
  CPU engines is 7--9\,\% slower on the GPU than N-I-J-K (\code{results/diamond\_engine\_cg19\_2026-11-04.txt}); 1\,024
  threads per tile, the pipeline maximum and the first default, cost 44\,\% on the notch tutorial and 21\,\% on the
  SMA launch (\code{results/diamond\_engine\_cg21\_2026-11-04.txt}). A one-dispatch fused kernel was not built: it
  needs a second copy of both fields, and $K=4$ already runs above the bound such a kernel has. On 2026-11-11 (cg31,
  Section~\ref{sec:cg31}), moving the Pre stages of Mur and of the conducting sheets into the half step before,
  beside the cell update, saved a barrier per half step with identical fields but ran 2--4\,\% slower, and was
  dropped.
\end{itemize}

\section{Discussion and future work}\label{sec:future}
Three lessons generalise. First, measure the stock code on the real models before believing a roofline: the stock
engine was not bandwidth-bound on the models that mattered, it was synchronisation- and instruction-bound, and the
largest single gains (barriers, inner loop, the allocation-placement valley, the dispersive hooks) were not about bytes
at all. Second, bit identity is a practical acceptance test for FDTD engine work: it is cheap to check, it forbids
the reorderings that would otherwise creep in, it caught one real bug (the hook plan, wrong in two snapshots), and it
confirmed the Debye bug that the review found by reading the code. Third, on a unified-memory laptop the GPU is
another client of the same bus, and it pays only once it moves fewer bytes per cell update than the CPU can. The
unfused GPU engines measured here (the base's and the Swift/Metal prototype) were at best about as fast as the fused
CPU sweep. Fusing alone is not enough: the
fused, unblocked Metal sweep ($K=1$ of the diamond engine) runs the two large models at 1\,730--1\,840\,\MC,
1.25--1.30$\times$ the unfused GPU engine of the fork's base but 0.58--0.65$\times$ of $K=4$ and 67--75\,\% of the
2\,470--2\,590\,\MC\ bound of Section~\ref{sec:ceiling}, 85--90\,GB/s at 49\,B per cell, because its per-step
structure (four dispatches, the tile edges, a threadgroup barrier between the two half steps) costs more than the
bytes it saves (paired A/B, 4 rounds, 2.5--4 cores of other load; \code{results/fused\_sweep\_k1\_2026-11-04.txt}).
Blocked over time steps the GPU runs past that bound, 2\,820--2\,990\,\MC\ at $K=4$ ($K=2$, at 2\,120--2\,410, stays
below it). Paired with the CPU engine of the same build at its best, 8 threads (Table~\ref{tab:current}: cg29, 4
rounds, other processes on a median of 0.5--0.8 cores), the Metal diamond engine ran patch x5 and notch x3
1.47$\times$ and 1.46$\times$ as fast and the GCPW line 1.17$\times$; it was level on the SMA launch on its
1.071\,M-cell mesh (1.00$\times$), and on the two tutorials the CPU engine was ahead (the GPU at 0.97$\times$ on the
notch tutorial and 0.93$\times$ on the patch tutorial). At cg19 the CPU engine was ahead on a $61^3$ Mur box too
(0.23\,M cells, 2\,281 against 2\,099\,\MC\ at 4 threads, rough numbers under a load of 2--7;
\code{results/diamond\_engine\_cg19\_2026-11-04.txt}): the GPU runs below its own rate on a model that does not fill it. In a size sweep of three
models (Figure~\ref{fig:crossover}) it led in every round from 1.7--2.9\,M cells on, by 1.25--1.52$\times$ above
3\,M; on smaller grids it waits for the CPU between short calls (18--27\,\% of the loop on the patch tutorial at
cg29) or runs at a lower rate of its own (the notch tutorial), and the CPU engine is at its fastest between 0.6 and 1.9\,M cells; below 1.7--2.9\,M
cells the two are within 0.89--1.23$\times$ of each other (Section~\ref{sec:gpu}). The build of 2026-11-11, cg31,
raised the Metal engine's rates on the six models by 1.07--1.53$\times$ (three dispatches per block, UPML and
dispersive stages on z-vectors or on the three components of a cell, seven tilings); in its paired rounds the Metal
engine is ahead of the CPU engine at its best on five of them, by 1.08--1.59$\times$, and level on the patch tutorial
(Section~\ref{sec:cg31}, Table~\ref{tab:cg31}). Once both engines are near what they can draw from
memory, the cheapest remaining wins are in the model. Still open: fp16 field storage (Section~\ref{sec:rejected}); larger Apple GPUs, on which luikore's
diamond engine (PR \#228~\cite{luikore228}) reports several times the CPU (Section~\ref{sec:positioning}; on this M5
it refused the six models, which all have Mur faces, and ran the PEC variants of five of them, speed-only cases, at
0.39--2.00$\times$ cg18, \code{results/others\_2026-11-04/README.md}; on its own benchmark model cg31's Metal engine
was level with it at 3.0\,M cells and 1.17$\times$ as fast at 17\,M, Table~\ref{tab:threebuilds}); and spreading one
model over several machines (Section~\ref{sec:cluster}).

On the engine side, a second read-only optimisation review, made on 2026-11-02 after the work up to then, ranked what
was left (Table~\ref{tab:review2}); its gains were estimates. Its CPU items were built and measured for cg29 on
2026-11-05 (Section~\ref{sec:cg29}): items 1, 2 and 5 are kept, item 9 in another form and item 4 in part; items 3,
6 and 7 gained nothing outside the noise and were removed (Section~\ref{sec:rejected}); item 10 was not built and item
8 was skipped. Against the estimates, item 1 gained 12.5--18.3\,\% on the SMA launch (estimate 15--25\,\%), item 2
3.1--4.1\,\% on the GCPW line (8--12\,\%), and item 5 more than expected on the two boards at 8 threads (28.6 and
10.9\,\% against 0--3 and 3--8\,\%). The cost of the nf2ff dumps on the smallest model remains (14\,\% at snapshot
cg8, project log; on cg18 at 4 threads 1\,845.9\,\MC\ with them against 2\,198.1 without,
\code{results/patch\_s1\_n10500\_base.txt}). The one item that review missed is done in cg29: the operator index,
allocated per cell up to cg28 although only one entry per z-vector is used, is now allocated per z-vector, which saves
up to 3\,B per cell of DRAM traffic on meshes of up to 42 z lines (bit-identical, since only the layout changes):
patch x5 1.119$\times$ at 4 threads and 1.101$\times$ at 8, notch x3 1.042$\times$ and 1.052$\times$, the other four
models 1.005--1.021$\times$ (paired A/B in one build against \code{OPENEMS\_OP\_INDEX\_PER\_CELL=1}, other processes
on 0.3--1.4 cores; \code{results/cpu\_items\_cg29\_2026-11-05.txt}, section 1).

\begin{table}[htbp]\centering\small
\setlength{\tabcolsep}{4pt}
\begin{tabular}{@{}>{\raggedright\arraybackslash}p{0.25\textwidth}>{\raggedright\arraybackslash}p{0.16\textwidth}>{\raggedright\arraybackslash}p{0.29\textwidth}lc@{}}
\toprule
idea & expected gain GCPW / SMA / big models & measured on 2026-11-05 (cg29) & effort (h) & bit-identical \\
\midrule
1. UPML y-vector path for PML boxes that are not full-z slabs (four consecutive y rows as one vector, lane gathers) & 0 / +15--25\,\% / 0 & kept (0019): SMA 1.183$\times$ (1 thr), 1.125$\times$ (4), 1.129$\times$ (8); 6 rounds, 0.4--0.6 cores & 8--12 & yes \\
2. NEON 4-position ADE update for conducting sheets (Lorentz line hooks) & +8--12\,\% / 0 / 0 & kept (0018): GCPW 1.041$\times$ (4 thr), 1.031$\times$ (8); 8 rounds, 0.3--0.5 cores & 5--8 & yes \\
3. Thread-count cache per model and a 4-candidate search & about 1.2\,s per run / +1\,\% / 0 & dropped: at real run lengths GCPW 1.021$\times$ [0.893--1.044], SMA 1.002$\times$; 5 rounds, 0.5--0.8 cores & 3 & yes \\
4. Stable P/E assignment by per-thread QoS, longer spin, faster slab balancing & 0 / 0--8\,\% at 8 thr / 0--8\,\% at 8 thr & dropped: at 8 thr QoS split 0.949--1.357$\times$, spin 200\,$\mu$s 0.907--1.122$\times$, 5\,ms balancing 0.991--1.008$\times$; 5 rounds, 0.3--0.5 cores; kept in part: no flag spin above the P-core count (0027) & 2--3 & yes \\
5. Neighbour flags instead of the two global barriers & 0--3\,\% / +3--8\,\% / +4--10\,\% at 8 thr & kept (0025): 1.015--1.055$\times$ (4 thr), 1.073--1.286$\times$ (8; GCPW 1.286$\times$, SMA 1.109$\times$, big models 1.077--1.102$\times$); 5 rounds, 0.3--0.6 cores & 10--16 & yes \\
6. Mur z-face vectors (lane gathers over y) & +3--8\,\% / 0 / 0 & dropped: 0.982--1.012$\times$ (GCPW, patch tutorial, patch x5); 6 rounds, 0.4--2 cores & 4 & yes \\
7. Per-line active masks for the line hooks & +2--4\,\% / +1--2\,\% / 0 & dropped: 0.960--1.020$\times$ (GCPW 1.005$\times$, 1.009$\times$); 5--8 rounds & 2 & yes \\
8. The base's Metal engine: one dispatch per hook for all Mur faces, skip empty hooks & GPU path +30--50\,\% on GCPW-class models (estimate) / +5--10\,\% / 0 & skipped: the Metal diamond engine (Section~\ref{sec:gpu}) replaces the base's GPU engine & 8 & GPU only \\
9. Split packed coefficient tables (geometry part 48\,B + a small material id) & 0 / +3--6\,\% / +1--3\,\% & split form dropped (SMA 0.952--0.999$\times$, GCPW 0.933--0.954$\times$); kept as geometry tables (0023): SMA 1.067$\times$ (8 thr), 1.033$\times$ (auto), 1.025$\times$ (4), 0.992$\times$ (1); 8 rounds, 0.3--0.4 cores & 3 & yes \\
10. Spin barriers for the start/stop of each processing interval & tutorials +3--8\,\% / 0 / 0 & not built: two wake-ups of 20--40\,$\mu$s per call of 50--300\,ms, under 0.1\,\%; patch tutorial 1--2\,\% at most (estimate) & 1--2 & yes \\
\bottomrule
\end{tabular}
\caption{Second optimisation review (read-only, 2026-11-02): the gain it expected on the GCPW line, the SMA launch and
the two large models (estimates, made before any item was built), effort in hours, whether the result stays
bit-identical, and the outcome measured on 2026-11-05 for snapshot cg29 (Section~\ref{sec:cg29}; paired interleaved
A/B in one build, median of the per-round ratios, rounds and the other processes' load where recorded, the SMA launch
on its 0.871\,M-cell mesh; \code{results/cpu\_items\_cg29\_2026-11-05.txt}, sections 1 and 2).}\label{tab:review2}
\end{table}

The review also located why threads still wait at the barriers: the fused sweep runs static slabs that are rebalanced
only between processing intervals, while macOS migrates threads between the clusters on millisecond timescales
(per-x-line stealing is used only by the unfused plan). cg29 measured the wait (at 8 threads 24--27\,\% of the time at
the two barriers, 10--22\,\% on the neighbour flags that replace them; Section~\ref{sec:cg29}). The review also
separates the two boards, by its own estimates rather than by measurement: the GCPW board looks instruction-bound in
the P-cluster cache (its conducting-sheet update about 20--25\,\% of a step), while the SMA launch looks bus-bound,
with its scalar PML path costing about 5.7 plain cells per PML cell. One part of that is now measured: with every
z-vector on one operator (\code{OPENEMS\_DBG\_OPMASK=1} in a development build, wrong fields, timing only) the SMA
launch ran 1.140$\times$, 1.128$\times$ and 1.076$\times$ faster at 4, 8 and 1 threads, the GCPW line 1.086$\times$
and 1.027$\times$ at 4 and 8, and the two large models within the noise, so only the SMA launch streams its
coefficient tables from memory (\code{results/cpu\_items\_cg29\_2026-11-05.txt}, section 2).

cg31 (Section~\ref{sec:cg31}) worked on the Metal engine's per-step structure and on its extension stages: three
dispatches per block instead of four, UPML and dispersive stages that take a z-vector or the three components of a
cell at a time instead of single floats, Mur faces of one axis behind one barrier, seven tilings instead of four.
Three items are open after it. On the CPU, temporal blocking with one cache-sized tile for all threads, as upstream
master's \code{cpu} engine does it: the blocking is worth 1.15--1.49$\times$ to that engine on this machine, and on a
17\,M-cell grid it is ahead of the fork's CPU engine (Section~\ref{sec:threebuilds}), while the tests of
Section~\ref{sec:rejected} gave every thread a sub-grid of its own. In the Metal kernel, the cost of an extension
that is compiled in: in timing-only builds of two snapshots the GCPW line ran 1.26--1.35$\times$ as fast without its
Mur faces and 1.22--1.32$\times$ without its sheets, but skipping one of their stages at run time (on the later of the
two) gave 1.02--1.07$\times$, and what in the GPU's execution makes it so was not found. The third is a shorter excitation pulse (next paragraph).

Open on the modelling side: the same mesh-floor rule for the other builders; a convergence pair (coarse/fine) as part
of every study; the SMA launch's top PML, eight cells over the whole x-y face although only the coax wave has to be
absorbed there (a thinner one was not validated on the full-size board); a baseband excitation pulse for boards on
which every conductor has a DC path (on the two boards it cut the steps 1.52$\times$ and 1.94$\times$ with $S_{21}$
within 0.005\,dB of the full run; a board with a conductor without a DC path, or with a series capacitor, where the
signals may not settle, was not tested; Section~\ref{sec:pulse}); and the time step. openEMS's default step
(\code{CalcTimestep\_Var3}) is a per-node bound that is exact on a uniform mesh and conservative at grading steps: on
the GCPW line the largest eigenvalue of the lossless update operator puts the leapfrog limit of the bulk update 1.26
times above it (computed for this report from the operator exported by cg18,
\code{results/timestep\_bound\_gcpw\_2026-11-03.txt}). A larger step was not adopted: the results would not be
bit-identical, every mesh needs an eigenvalue solve, and stability with the Mur, PML, lumped and dispersive
extensions, which that bound leaves out, is not proven.

Upstreaming: the CPU engine changes are platform-independent apart from a few guarded details: thread QoS and the
P-core count that starts the thread search and, since cg29, ends the spin of the neighbour flags (macOS; other systems
start at half the available hardware threads, and their flag waits spin like the barriers), and the opt-in
flush-to-zero (arm64); the separable nf2ff path needs a CBLAS, and only Apple's Accelerate is wired in. The Metal
diamond engine is Apple-only: it is built with the base's Metal backend (\code{ENABLE\_METAL}, on by default on Apple)
and runs only when \code{--engine=metal-diamond} asks for it. The allocator stagger is not guarded: it applies on
every platform, although the valley it removes was seen only with the macOS allocator, and on the M5 it costs nothing
measurable outside the valley (Section~\ref{sec:valley}). The
fork is based on commit 7bb294f (2026-10-12) of upstream's \code{gpu-engine} branch. Upstream has since merged a
reworked GPU engine into master (c3a956b, 2026-10-17), deleted the branch, and replaced the operator and engine
classes the patches change (c0816a2, 2026-11-02), so a submission means porting each item and verifying it again
against a stock build of master. Upstream has also put speed pull requests on hold until a refactoring of the
operator classes is done~\cite{thliebig495}, and several items here overlap pull requests already in that queue
(\#228, whose tile schedule the Metal diamond engine follows; \#231; \#244, whose UPML part cg29 takes; \#245, whose
strict-order check cg29 follows; Section~\ref{sec:positioning}), so a submission would start from what is still
missing after them. Upstream's \code{AGENTS.md} also asks that a pull request claiming a speed-up include
the \code{benchmark/bench.py compare} table of master against the branch, run interleaved on one machine, and that no
dependency be added. Every case of that benchmark uses a sinusoidal source, which turns on the steady-state detection;
up to cg28 that extension had no line hooks, so the fork ran those cases on the unfused plan. Since cg29 its line
hooks and \code{TS\_LOCAL} put them on the fused sweep and the neighbour flags (fused against unfused: pml
1.36$\times$ at 4 threads, inhomogeneous 1.29$\times$ and 1.34$\times$ at 4 and 8, the dispersive case, with Debye
line hooks too, 6.5$\times$ and 10.0$\times$; Section~\ref{sec:cg29}). The tool was run on master and on the fork
on 2026-11-10 (Section~\ref{sec:common}, Table~\ref{tab:upstream}), without its probes column, which needs master's
Python module: on the \code{multithreaded} engine the fork is 1.23--10.0$\times$ master, and master's new \code{cpu}
engine is level with or ahead of the fork's CPU engine at 8\,M cells (on luikore's benchmark model, run interleaved
on 2026-11-11, cg31's CPU engine was 1.27$\times$ master's \code{cpu} engine at 3.0\,M cells and 0.86$\times$ at
17\,M; Table~\ref{tab:threebuilds}), so a submission of the CPU items would have
to show a gain over that engine, not over \code{multithreaded}. The Metal diamond engine does not run the steady-state detection on the GPU: with it, every step runs
the single-threaded CPU loops of the compressed engine (Section~\ref{sec:gpu}). A submission would also follow
upstream's AI policy~\cite{aipolicy} (a \code{Generated-by:} trailer, a sign-off, and a human who has read every
line).

\section{Reproducibility and artefacts}\label{sec:repro}
fdtdKit, the subject of this report, is the FDTD kit of EMSKit, a family of electromagnetic solvers for Apple
silicon; the family's other kits (femKit for finite elements, momKit for the method of moments and sbrKit for shooting
and bouncing rays) are developed separately and are not covered here. Until 2026-11-05 fdtdKit was called emsKit, and
the repository and its folder keep that name, as do the result files, the patches and the dated entries of
\code{CHANGELOG.md}.
The repository that accompanies this report, \repourl{} (archived at Zenodo, DOI
\href{https://doi.org/10.5281/zenodo.24118865}{10.5281/zenodo.24118865}), contains:
\begin{itemize}
\item \code{patches/openems/}: the fork as \code{git am} patches 0001--0039 on commit 7bb294f of upstream's
  \code{gpu-engine} branch (\code{BASE\_origin\_gpu-engine.txt}; upstream has since deleted the branch, so
  \code{BUILD.md} fetches the commit by its hash), and one cumulative diff that gives the same tree (dc4a352).
  0001--0006 are
  cg18, the build installed on the author's Mac from 2026-11-04 to 2026-11-06; 0007--0016 add the Metal diamond
  engine, one patch per snapshot cg19--cg28; 0017--0027 are the CPU items of cg29 (Section~\ref{sec:cg29}), so
  0001--0027 give cg29, the build installed there from 2026-11-06 to 2026-11-11; 0028--0039 are the twelve commits
  of 2026-11-11 (72a6492 to ff2538a, Section~\ref{sec:cg31}), so 0001--0039 give cg31, the build installed there
  since; the result files carry these snapshot names. All 39 patch files were made again on 2026-11-11; in
  0001--0027 only the numbering in the subject line changed.
  The fork's commits up to cg29 were rewritten for publication on 2026-11-07, with the author's public identity,
  the trailers
  of upstream's AI policy and the same trees, and pushed that day; the result files made before then and the
  entries of \code{CHANGELOG.md} dated before then give their earlier hashes, which \code{CHANGELOG.md} maps to the
  published ones. The version strings that the libraries of those snapshots print keep the hashes at which their
  build folders were configured, from before the rewrite; the fresh build of 2026-11-07 (pub29) prints the hash of
  its own \code{git am} commit, v0.37.0-45-g1d091c7, and the installed cg31 that of the fork's head,
  v0.37.0-57-gff2538a.
  \code{patches/csxcad/}: one patch on v0.7.0 (+9 lines), which the openEMS patches need. The changes made in this
  project to the author's simulation toolkit (rfsim) are not published, since the toolkit itself is private;
  Section~\ref{sec:modelling} and \code{CHANGELOG.md} describe them.
\item \code{bench/}: \code{bench.py}, \code{runner.py}, \code{models.py}, \code{check.sh}, \code{step1.sh}, and since
  2026-11-04 \code{usage.py} (the other processes' CPU and GPU use during a run, which every benchmark line now
  carries) and \code{abtest.py} (the paired interleaved A/B of Section~\ref{sec:cg29}). The four
  tutorial models run on any openEMS install; the two boards import the author's rfsim toolkit, which is private:
  neither it nor the changes made to it in this project are published, so they run only where the toolkit is installed
  (\code{RFSIM\_PATH}). \code{bench/validation/}: 15 checks, 12 against closed-form or exact references (dispersive
  slabs, a rectangular and a cylindrical cavity, a WR-90 guide, a microstrip line) and three that compare the engine
  under test with the basic engine of the same build on the dispersive slabs. Its runs are in
  \code{results/validation\_<tag>.md} (and \code{.json}). On the M5 cg18 passes 15 of 15 with the
  multithreaded engine at 4 threads (\code{cg18\_mt4}) and with the default engine and thread count
  (\code{cg18\_fastest}), and so do cg29 (\code{cg29\_mt4}, \code{cg29\_fastest}; \code{cg29pre\_*} are an earlier
  build of it), the fresh build of 2026-11-07 of the patches published then, 0001--0027 (\code{pub29\_mt4},
  \code{pub29\_fastest}), and stock v0.37.0 with its defaults (\code{stock-v0.37.0\_macos}); the Metal diamond engine passes 15
  of 15 at cg23 and cg27 (\code{cg23\_metal}, \code{cg27\_metal}; \code{VALIDATION\_ENGINE} selects the engine, the
  cylindrical cavity runs the cylinder engine whatever it says, and the dispersive slabs on the GPU equal the basic
  engine of the same build); snapshot cg17
  (patches 0001--0005) passes 15 of 15 with its defaults and with the multithreaded engine at 8 threads (\code{cg17},
  \code{cg17\_mt8}, two runs not recorded in this project's log). Earlier runs:
  stock v0.37.0 on a Linux x86\_64 virtual machine passed 12 of 12 (\code{stock-v0.37.0\_linux-sandbox}, the suite's
  earlier version without the three engine comparisons). Two negative controls fail as they should: a stock build
  with the Debye engine hooks disabled, on that machine, passes 2 of 3 (\code{negative-control\_debye-hooks-disabled});
  on the M5 a build of patches 0001--0004 (snapshot cg16, which has the Debye bug) with the multithreaded engine at
  2 threads passes 4 of 6, and both Debye checks fail (\code{negative-control\_patches-0001-0004\_mt2}).
\item \code{results/}: the benchmark tables, one file per model, step count and snapshot, with the load figure of
  every run, among them the full checks of cg18 (\code{*\_cg18.txt}) and of a fresh build of the patches published on
  2026-11-07, 0001--0027 (\code{*\_pub29.txt}), and the comparison of stock, the fork's base and cg18 (\code{*\_base.txt});
  \code{summary.json} (best speed per model and snapshot; the cg17 tables and every table after cg18, those of cg19,
  cg21, cg22, cg29, cg29pre and pub29 and the \code{*gpu.txt} tables of cg20 to cg29 and pub29, are correctness
  checks with unpaired runs, many of them on a loaded machine, so their speeds do not compare; for a
  \code{*\_cg2Ngpu.txt} or \code{*\_pub29gpu.txt} table the entry is its multithreaded row with
  \code{OPENEMS\_FTZ=1}); the mesh and cache-valley sweeps and the valley's
  placement test; the end-to-end runs of the two boards (\code{rfsim\_e2e\_2026-11-03.txt}, and on the GPU
  \code{rfsim\_e2e\_gpu\_2026-11-04.txt}); the kernel comparisons (GCC against clang, instructions per loop
  iteration); the correctness checks of cg18; the validation runs; the memory and time-step tables; the initial
  review. From 2026-11-04 and 2026-11-05: the Metal diamond engine's snapshots
  (\code{diamond\_engine\_cg19\_2026-11-04.txt} to \code{diamond\_engine\_cg28\_tuner\_2026-11-04.txt}, the Swift
  diamond mode \code{diamond\_swift\_2026-11-04.txt}, the thread sweep \code{diamond\_sweep\_cg21\_2026-11-04.txt}, the
  \code{diamond\_check} logs of cg21--cg23, the reduced CPU checks \code{*\_cg19.txt}, \code{*\_cg21.txt} and
  \code{*\_cg22.txt}, and the GPU-against-CPU tables \code{*\_cg}$N$\code{gpu.txt} of cg20--cg23 and cg27--cg29); the
  $K=1$ measurement (\code{fused\_sweep\_k1\_2026-11-04.txt}); the runs of upstream master
  and of PR \#228 (\code{others\_2026-11-04/}, with its \code{README.md}); the CPU items of cg29
  (\code{cpu\_items\_cg29\_2026-11-05.txt}, with the identity check tables \code{*\_cg29.txt}, whose speeds are the
  faster of two unpaired runs under the check's load, and \code{*\_cg29pre.txt} of an earlier build); the GPU against
  the CPU by grid size on cg29 (\code{gpu\_crossover\_cg29\_2026-11-05.txt}, Figure~\ref{fig:crossover}); and the
  tables of cg29:
  stock, cg29 on the CPU and the Metal diamond engine side by side (\code{speed\_cg29\_2026-11-05.txt},
  Table~\ref{tab:current}) and the end-to-end runs of cg29 on the CPU and the GPU
  (\code{rfsim\_e2e\_cg29\_2026-11-05.txt}, Table~\ref{tab:e2ecg29}), from which \code{paper/figures/tables.py}
  writes those two tables and the speed table of \code{README.md}. From 2026-11-10 and 2026-11-11: the common
  benchmarks (\code{common\_benchmarks\_2026-11-10.txt}, Section~\ref{sec:common}) and the records of the current
  build, cg31 (Section~\ref{sec:cg31}): its Metal engine against cg29's (\code{metal\_engine\_cg31\_2026-11-11.txt}),
  the three builds on luikore's model (\code{engines\_on\_luikore\_model\_2026-11-11.txt}; with cg29 the night before,
  \code{luikore\_compare\_cg29\_2026-11-11.txt}) and the pulse test (\code{pulse\_test\_2026-11-11.txt}).
\item \code{scripts/} (one line per script in \code{scripts/README.md}): the standalone kernels (\code{kb\_idx.cpp},
  \code{kb\_tb.cpp}) with the GCC-against-clang comparison (\code{gcc\_vs\_clang.sh}) and the instruction count per
  loop iteration (\code{loop\_insns.py} with the program-counter sampler \code{pc\_sampler.cpp}), the cache-valley
  box, sweep and placement test (\code{place.c}, \code{placement.sh}), the cylindrical-multigrid check, the
  setup/nf2ff phase timers, the probe-end and z-padding checks, the mesh sweep, the end-to-end runs
  (\code{scripts/e2e/}; \code{E2E\_ENGINE=metal-diamond} in \code{run\_one.py} runs them on the GPU), the memory and
  time-step checks, the HDF5 compare, the Debye regression check of patch 0005 (\code{scripts/checks/debye\_check.sh}),
  the Metal diamond engine's bit-identity check against the multithreaded engine with \code{OPENEMS\_FTZ=1}
  (\code{scripts/checks/diamond\_check.sh} with \code{diamond\_box.py} and \code{diamond\_compare.py}), the
  GPU-against-CPU size sweep (\code{scripts/gpu\_crossover/}: \code{sweep.sh}, \code{tiles.py}, \code{cells.py},
  \code{analyze.py}, and \code{mkresult.py} with \code{header.txt}), the builds and runs of the other engines
  (\code{scripts/others/}: \code{run\_others.sh}, \code{bench\_ext.py}, \code{xcompare.py}), and from 2026-11-10 and
  2026-11-11 the common benchmarks (\code{scripts/common\_bench/}), the three builds on luikore's model
  (\code{scripts/luikore\_compare/}), the closed cylinder (\code{scripts/cylinder\_closed/}) and the pulse test
  (\code{scripts/pulse/}).
\item \code{swift\_engine/}: the Metal kernel, with its diamond mode (\code{--diamond K}), its model builder and field
  compare.
\item \code{Package.swift}, \code{Sources/} and \code{Tests/}: fdtdKit's command line \code{emskit fdtd} (version
  0.1.0), a subcommand of the \code{emskit} tool of the EMSKit Swift package. It runs the tools above unchanged
  (\code{bench.py}, \code{abtest.py}, \code{check.sh}, the validation suite, the Debye and diamond checks, the
  end-to-end runs, or any openEMS Python script) and prints one JSON document per call; every speed and run time it
  prints carries the other processes' CPU and GPU use during that run, it installs nothing, and it refuses any
  location it would write inside the installed openEMS prefix or the snapshot folder. Its plan, with the log of its
  builds, tests and acceptance runs, is \code{Docs/fdtdKit/PLAN.md}. Written on 2026-11-05 and 2026-11-06 and
  corrected on 2026-11-09 and 2026-11-10 after two code reviews (Section~\ref{sec:disclosure}), it was last built
  and tested on 2026-11-10 with Swift~6.4.2 on macOS~27.0.1: its 244 unit tests pass (258 tests in 28 suites with
  the 14 smoke runs against the real tools). The smoke runs passed last earlier that day, before the day's later
  corrections, as did the runs on cg29: \code{check --lite} gave results identical to stock on the six models,
  \code{diamond-check} 72 of 72, and the GCPW end-to-end run finished (34\,300 steps). The checks of the installed cg29
  (Section~\ref{sec:method}) were run with it.
\item \code{paper/}: this report, its figures and the scripts that draw them, and \code{code\_notes.md}, the
  description of every code change of patches 0001--0006 written from the diffs (Section~\ref{sec:engine}; the later
  patches are described by their commit messages, \code{CHANGELOG.md} and Sections~\ref{sec:gpu},
  \ref{sec:cg29} and~\ref{sec:cg31}). Six figure scripts read
  \code{results/}: the cache-valley figure reads its two sweep files, the crossover figure
  \code{gpu\_crossover\_cg29\_2026-11-05.txt}, and the speed, thread, progress and bandwidth
  figures read \code{summary.json} (written by \code{figures/parse\_results.py}); the bandwidth figure adds four bars
  from the project log. The mesh-sweep and run-time figures hold their values in the script, copied from
  \code{results/sma\_mesh\_sweep\_2026-11-01.txt} and \code{results/rfsim\_e2e\_2026-11-03.txt}.
\end{itemize}
\code{BUILD.md} gives the build, install, verification and revert steps. Licences: openEMS is under the GPL and CSXCAD
under the LGPL, both version 3 or later. The openEMS patches are under the GPL like the files they change, the CSXCAD
patch under the LGPL like its files (the licence text is in \code{patches/csxcad/}); files copied from or based on
openEMS keep their upstream copyright and licence notices; the harness, the scripts (the figure scripts in
\code{paper/figures/} included), \code{paper/code\_notes.md}, \code{swift\_engine/} and the result files of this
repository are under GPL-3.0-or-later (\code{LICENSE}); the text and figures of this report are under CC BY 4.0
(\code{paper/LICENSE-CC-BY-4.0.txt}). The Swift package of the command line (\code{Package.swift}, \code{Sources/},
\code{Tests/}) is under GPL-3.0-or-later: its Swift files and its Python launcher (\code{fdtdkit\_run.py}) each carry
\code{SPDX-License-Identifier: GPL-3.0-or-later}, and \code{LICENSE} holds the licence text.

\section{Disclosure}\label{sec:disclosure}
The code changes, the benchmark harness, the command line, the measurements and most of this text were produced by AI
assistants (Claude, Anthropic, operated through Claude Code) over fourteen days (2026-10-28 to 2026-11-10), directed by
the author through a chat interface; the models are those recorded in the sessions' transcripts. In the main session:
Claude Opus 5.5 on 2026-10-28 until the evening (the benchmark harness, the stock baseline and the engine work of
patch 0001, snapshots A to cg8); Claude Fable 5.1 from then to the early morning of 2026-11-03 (it made the commits
of patches 0001--0005, so their trailers name it, and the trailer of 0001 also names Claude Opus 5.5, which wrote
that patch's code; it made the engine work of snapshots cg9 to cg17, the CSXCAD patch and the toolkit changes up to
2026-11-02, and most of the first version of this report);
and Claude Opus 5.5 from then on: the final review of 2026-11-03 with patch 0006 (below), on 2026-11-04 the Metal
diamond engine (patches 0007--0016) with its checks and the $K=1$ measurement, and on 2026-11-05 the CPU items of cg29
(patches 0017--0027), the measurements of Tables~\ref{tab:current} and~\ref{tab:e2ecg29} and the documentation pass
that brought this report up to date with them, whose fact sheets and edits were made by Claude Opus 5.5 agents in
workflows of the main session. The main session also ran the GPU-against-CPU size sweep of
Figure~\ref{fig:crossover} on 2026-11-05; five checkers and a skeptic, Claude Opus 5.5 agents in a read-only
workflow, then recomputed its rows from the raw runs and checked its text. At the author's request, Claude Opus 5.5
agents in workflows of the main session renamed the documents to fdtdKit on 2026-11-05 and made the command line
\code{emskit fdtd} on 2026-11-05 and 2026-11-06: its specification, code and tests, its build, its reviews and its
acceptance runs. On 2026-11-06 the main session and Claude Opus 5.5 agents installed cg29 on the author's Mac,
checked it from the install and brought this report up to date with the command line and the install. Later that
day the main session and Claude Opus 5.5 agents in one of its workflows applied the author's publishing decisions to
the repository and to this report (the author line, the repository's address, the licence of the command line's
Swift package, and leaving the toolkit changes out).
On 2026-11-07, at the author's request, the main session rewrote the commits of the two fork branches for
publication (the author's public identity and the trailers of upstream's AI policy, with the same trees), compared
every rewritten commit with its old one and pushed the two branches to the author's GitHub forks; it then made the
patches again from the published commits and replaced the hashes that point to the fork's code in this report and
in the repository's files (Section~\ref{sec:repro}). Claude Opus 5.5 agents in two read-only workflows of the main
session checked this work: before the push the first rewrite and a first version of the changed files, after it the
pushed branches and the final files. The same day the main session built the published patches from a fresh clone
and ran the full check on that build (Section~\ref{sec:method}); Claude Opus 5.5 agents in a third read-only
workflow checked its record against the raw logs, and the changed files. Later that day Claude Opus 5.5 agents in a
fourth read-only workflow (six auditors, a skeptic and a critic) audited the copy of the repository made for
publication; on 2026-11-08 Claude Opus 5.5 agents in another workflow of the main session applied the fixes that the
audit proposed and the author's decisions of that day (among them the licence wording of Section~\ref{sec:repro}) to
the repository and to this report. The same day the author ran Claude Code's cloud code review
(\code{/code-review ultra}; the model of its reviewers is not shown) on 18 of the command line's 31 source files
(7\,867 of the 11\,634 lines of the reviewed version: the files that start programs, choose the libraries, write or
delete files and decide pass or fail). It reported one bug (a binding without package metadata was taken for one that failed to import) and
five minor issues. On 2026-11-09 a separate Claude Opus 5.5 session, opened by the author in the review's copy, wrote
the six fixes and four tests. Claude Opus 5.5 agents in a read-only workflow of the main session checked them; the main
session corrected what they found (one fix had dropped a fallback; the handling of an interruption in another needed
three small changes), added five tests and carried all of it into the repository.
From the evening of 2026-11-09 the main session ran on Claude Fable 5.2: its agents reviewed all 31 source files of
the command line in a read-only workflow (four reviewers, a refuting verifier for each, a critic), and on 2026-11-10
agents of a second workflow fixed the confirmed findings and reviewed the fixes, three times; the main session ran the
checks.
After that period, on 2026-11-11, the main session (Claude Fable 5.2) made the build cg31 (patches 0028--0039), its
measurements and Section~\ref{sec:cg31}; the search for the bottlenecks before it was the work of four read-only
reviewers, each followed by a checker that tried to refute its findings (agents of the same model in a workflow of
the main session), and the day before one read-only checker had recomputed the numbers of the common benchmarks of
Section~\ref{sec:common}.
A separate Claude Opus 5.5 session in the
repository's folder wrote the paragraph on concurrent work and measured two of the engines it lists (below); another
Claude Opus 5.5 session wrote the work plan of 2026-11-04 (a working note that is not part of the published
repository), which the author pasted into the main session. The multi-agent reviews and the research of 2026-11-04
and 2026-11-05 (a performance review of the
diamond kernel, a review of its tiling tuner, the reading of the other contributors' pull requests and fork) ran as
read-only workflows of the main session with Claude Opus 5.5 agents. Since 2026-11-07 the fork's commits carry
the author's public git identity and name the model in a \code{Generated-by:} trailer above a
\code{Signed-off-by:} line. The goal, the acceptance rule (bit identity,
measured on their models), the plans and the choice of what to pursue next and what to reject were the author's, and so
were the decisions recorded with the work, among them the toolkit defaults of 2026-11-04 and the rule of the same day
that every benchmark line carries the other processes' CPU and GPU use. The mesh-floor recommendation of
Section~\ref{sec:meshfloor} (0.05\,mm with a hold of 0.5) came from the assistants; it stands on the numbers
re-measured on 2026-11-03, and the author adopted it as the toolkit's default on 2026-11-04 (toolkit commit 4a7e20e;
on 2026-11-03 the toolkit still had the earlier defaults, 0.02\,mm and a hold of 1). Following arXiv's policy the
assistants are not listed as authors. Two read-only reviews of the branch,
besides the initial review of Section~\ref{sec:baseline}, by fresh instances of Claude Fable 5.1 (one for
correctness, one for further optimisation) found no confirmed bug and located the conducting-sheet problem of
Section~\ref{sec:engine}. A second round of four such reviews (fork correctness, further optimisation, the rfsim
changes, and this report's numbers against the result files) was run before publication; the correctness
review found a real wrong-result bug (a Debye material losing its dispersion in the multithreaded engine, fixed in
patch 0005, Section~\ref{sec:reviewfixes}); the findings of all four are folded into Sections~\ref{sec:engine},
\ref{sec:modelling} and~\ref{sec:future}, and the corrections they forced are listed in the repository's
\code{CHANGELOG.md}. Section~\ref{sec:positioning} (with the abstract's sentence on clusters), \code{ROADMAP.md} and the
validation suite \code{bench/validation} with its two runs on the Linux virtual machine were added on 2026-11-03 by an
AI assistant session that the project log does not record; its model is not recorded. The two runs of snapshot cg17
on the M5 (\code{results/validation\_cg17*}) came later that day from a Claude Opus 5.5 session in the Claude desktop
app, also outside the project log.
The final review of 2026-11-03 was run by Claude Opus 5.5 (Anthropic): 30 reviewers, every finding checked by two
verifiers, the nits by one. Its fixes are in patch 0006 and in these documents, and a large part of the present text
was rewritten or added in that round. In the same round the session that coordinated the review (Claude Opus 5.5)
re-measured the end-to-end runs, on stock v0.37.0 and on cg18 (patches 0001--0006), and the SMA mesh sweep on cg18
(\code{results/rfsim\_e2e\_2026-11-03.txt}): Tables~\ref{tab:e2e},
\ref{tab:probeend} and~\ref{tab:mesh} and Figures~\ref{fig:runtime} and~\ref{fig:mesh} show those runs
(Table~\ref{tab:mesh} and Figure~\ref{fig:mesh} also keep one run of 2026-11-01 as history). The same session also
made the full check of cg18 against stock (\code{results/*\_cg18.txt}), the comparison of stock, the fork's base and
cg18 (\code{*\_base.txt}), the valley's placement test on cg18 and its sweep on stock, and the GCC/clang comparison in
the standalone kernel with the instruction counts of the inner loops. It also made the correctness checks of cg18 and
three validation runs on the M5 (cg18 twice, stock once), and re-ran the memory and time-step checks on cg18. The
validation suite's negative control on the
M5 came from the final review itself. The author has not yet read every line of the fork; this draft is marked
accordingly.
On 2026-11-04, at the author's request, Claude Opus 5.5 (Claude Code) added the paragraph on concurrent work of
Section~\ref{sec:positioning} with its references and narrowed the sentences it affects (the abstract, the
introduction, Sections~\ref{sec:ceiling}, \ref{sec:positioning}, \ref{sec:rejected}, \ref{sec:future} and this
one's upstreaming paragraph); it read the public pull requests, discussions and benchmark notes it cites, and when it
wrote the paragraph it had built or run none of that work. Later that day the same session built upstream master
(5fecbb5, with the Metal engine merged as c3a956b) and luikore's PR \#228 (55ce11b) and ran them on the six models next
to stock and cg18 (\code{results/others\_2026-11-04/}, Section~\ref{sec:positioning}).

Where the numbers come from: the benchmark tables (with the full checks of cg18, cg29 and the fresh build of
2026-11-07 and the comparison of stock, the fork's base and cg18), the end-to-end and probe-end runs
(Tables~\ref{tab:e2e}, \ref{tab:e2ecg29} and~\ref{tab:probeend}), the mesh sweep (Table~\ref{tab:mesh}), the valley
sweeps, the kernel comparisons, the correctness checks of cg18, the validation runs, and the measurements of
2026-11-04 and 2026-11-05 (the Metal diamond engine, the $K=1$ sweep, the other engines, the CPU items of cg29, the
GPU-against-CPU size sweep of Figure~\ref{fig:crossover} and the table of cg29 against stock, Table~\ref{tab:current})
and the records of 2026-11-10 and 2026-11-11 that Sections~\ref{sec:common} and~\ref{sec:cg31} name
are files in \code{results/}. So are the smaller checks made for this report during the final review, which are marked
where they appear: the memory figures of Section~\ref{sec:cluster}, the time-step bound of Section~\ref{sec:future}
and the placement test of Section~\ref{sec:valley} (all three on cg18), and the flush-to-zero check of
Section~\ref{sec:policy} (snapshot cg17); only the rerun of the Debye check with 2 threads
(Section~\ref{sec:reviewfixes}) is not. Two sets of checks of 2026-11-06 are not in \code{results/} either: the
command line's acceptance runs and the checks of the installed cg29 (Section~\ref{sec:method}), both logged in
\code{Docs/fdtdKit/PLAN.md}, section 16. A few figures of Section~\ref{sec:cg29} come from the commit messages of the
published patches and are marked so. Numbers that exist only in the project log, which is not published, are marked as
project-log numbers where they appear; among them, from before the final review, the setup and far-field timings, the
$201^3$-box figures and the standalone-kernel figures, and part of the list of Section~\ref{sec:rejected}. Numbers
that were computed or estimated rather than measured say so: the bytes per cell of the traffic model, the bounds of
Sections~\ref{sec:ceiling} and~\ref{sec:positioning} (the fused GPU bound now has the $K=1$ measurement beside it),
and the expected gains of Table~\ref{tab:review2}, which now stand beside their measured outcomes.
